\documentclass[11pt]{article}
\usepackage[margin=1in]{geometry}
\usepackage{amsmath,amssymb,amsfonts}
\usepackage{graphicx}
\usepackage{booktabs}
\usepackage{hyperref}
\usepackage{url}
\usepackage{enumitem}
\usepackage[numbers]{natbib}

\title{Quark: Contrastive Embeddings for Quantum Circuit Retrieval}
\author{[Sampark Bhol] \\ \texttt{contact@gmail.example}}
\date{\today}

\begin{document}
\maketitle

\begin{abstract}
We present Quark, a contrastive learning approach to embedding quantum circuits as fixed-size vectors for downstream search, deduplication, and clustering tasks. A small transformer encoder ($\sim$650K parameters) is trained with a symmetric InfoNCE objective on synthetic (anchor, positive, negative) triplets where positives are generated by composing equivalence-preserving rewrites (Hadamard cancellation, CNOT cancellation, gate commutation, $R_z$ angle splitting/merging). On in-distribution retrieval, Quark achieves Recall@10 = 1.000 in a 500-circuit pool, exceeding gate-count-vector cosine retrieval by 32 points. On a held-out rewrite the model never saw during training (\texttt{insert\_id}), Quark achieves Recall@10 = 0.935 versus 0.290 for the strongest baseline, providing evidence the model learns transferable structural invariances rather than rewrite-rule signatures. On QASMBench, where the task is to retrieve the transpiled lowering of an algorithm given its high-level form, Quark achieves Recall@10 = 0.345 versus 0.172 for gate-count-vector --- a real but small advantage in a strongly out-of-distribution setting. We release weights, code, and evaluation harness under Apache 2.0.
\end{abstract}

\section{Introduction}

The growth of quantum software repositories raises practical questions that classical ML treats as routine: how do you find duplicate or near-duplicate circuits, search a library by example, or cluster circuits by structure? Existing tooling treats circuits as objects to \emph{run} \cite{qiskit, pennylane}, not as data to \emph{index}. Exact equivalence checks (full unitary comparison) are exponential in the number of qubits and infeasible above $\sim$10 qubits. Syntactic identifiers (gate-sequence hashes, sorted gate lists) miss commutativity and rewrite-induced variation.

Quark addresses this gap with a learned embedding: a transformer encoder maps quantum circuits to unit vectors in $\mathbb{R}^{128}$ such that circuits known to be equivalent under unitary identities are placed close in cosine distance, while structurally unrelated circuits are placed far apart. Useful as building infrastructure for circuit deduplication CLI tools, search-by-example interfaces, and as a fast filter prior to exact unitary comparison.

\paragraph{Contributions.}
\begin{itemize}[topsep=2pt,itemsep=2pt]
    \item A circuit-as-data tokenization (gate type, qubit indices 1-indexed, parameter $\sin/\cos$ encoding) and a small transformer architecture for circuit embedding.
    \item An InfoNCE-trained encoder, evaluated against three principled baselines (gate-sequence MD5, sorted-gate MD5, gate-count cosine).
    \item A held-out-rewrite generalization test that distinguishes structural learning from rewrite-rule memorization.
    \item Real-corpus retrieval on QASMBench, including the algorithm-vs-transpiled-form retrieval task.
\end{itemize}

\section{Method}

\paragraph{Tokenization.} A circuit of $L$ gates is encoded as a sequence of $L$ tokens (capped at 128 with right-padding). Each token has four features: an integer gate ID over a 16-element vocabulary covering common one- and two-qubit gates plus two fall-through buckets for unrecognized gates; the index of the first qubit acted on (1-indexed, 0 = padding); the index of the second qubit (0 for one-qubit gates); and the cosine and sine of the rotation angle (defaulting to $\cos = 1, \sin = 0$ for non-parametric gates). The four features are mapped to a shared $d$-dimensional embedding by separate embedding tables (gate, q1, q2) and a linear projection (parameters), then summed.

\paragraph{Encoder.} A standard pre-norm transformer encoder ($d = 128$, 4 heads, 3 layers, GELU activation, learned position embeddings) reads the token sequence with a prepended \texttt{[CLS]} token. Padding positions are excluded via a key-padding mask. The final \texttt{[CLS]} hidden state is linearly projected and L2-normalized to produce the embedding $e \in S^{127} \subset \mathbb{R}^{128}$.

\paragraph{Training data.} Each training triplet $(a, p, n)$ is generated independently:
\begin{itemize}[topsep=2pt,itemsep=2pt]
    \item Anchor $a$: a random circuit with $n_q$ qubits ($n_q \sim U\{2,4\}$) and depth $d \sim U\{8,24\}$ where each gate is sampled from a distribution biased toward 2-qubit gates (CNOT) and Clifford operations.
    \item Positive $p$: $a$ after a chain of $k=4$ randomly chosen equivalence-preserving rewrites: $R_{HH}$ (insert/cancel a Hadamard pair), $R_{XX}$, $R_{CNOT}$, gate commutation on disjoint qubits, and $R_z(\theta) \rightarrow R_z(\theta/2) R_z(\theta/2)$ split/merge.
    \item Negative $n$: an unrelated random circuit drawn from the same distribution.
\end{itemize}
We verify that all rewrites preserve unitary equivalence using \texttt{qiskit.quantum\_info.Operator}, comparing matrices up to global phase.

\paragraph{Loss.} Symmetric InfoNCE \cite{infonce, clip} with temperature $\tau = 0.07$. Within a minibatch of $B$ embedding pairs $(e_a^i, e_p^i)$, the logits matrix $S_{ij} = e_a^i \cdot e_p^j / \tau$ produces a categorical loss:
$$\mathcal{L} = \frac{1}{2}\left[\text{CE}(S, I) + \text{CE}(S^\top, I)\right]$$
where $I$ is the identity assignment. In-batch negatives provide $B-1$ harder negatives per anchor than the explicit random negative used by triplet loss.

\paragraph{Optimization.} AdamW with $\beta_1 = 0.9, \beta_2 = 0.999$, weight decay $10^{-5}$, learning rate $3 \times 10^{-4}$ with cosine annealing, gradient clipping at norm 1.0, batch size 32, 12 epochs on 1500 generated triplets. Total training time: $\sim$5 minutes on CPU.

\section{Evaluation}

\paragraph{Setup.} For each evaluation, a query set of 100 random circuits is generated, and a library is constructed by including 2 known-equivalent variants (each a $k=4$ rewrite chain applied to its query) plus 298 unrelated random fillers, then shuffled. We report Recall@\{1, 5, 10\}: for each query, what fraction of its $k=2$ truth equivalents land in the top-$k$ retrieved results by cosine similarity.

\paragraph{Baselines.} (1) MD5 hash of the gate sequence; (2) MD5 hash of the alphabetically sorted gate-list; (3) cosine over a gate-count bag-of-words vector. These are the ones a reasonable engineer would try before reaching for ML.

\subsection{In-distribution retrieval}

\begin{table}[h]
\centering
\begin{tabular}{lccc}
\toprule
method & R@1 & R@5 & R@10 \\
\midrule
hash & 0.095 & 0.105 & 0.105 \\
sorted hash & 0.205 & 0.245 & 0.245 \\
count vec (cosine) & 0.305 & 0.555 & 0.675 \\
\textbf{quark (ours)} & \textbf{0.500} & \textbf{1.000} & \textbf{1.000} \\
\bottomrule
\end{tabular}
\caption{In-distribution retrieval, 100 queries / 500-circuit library, 2 truth equivalents each.}
\end{table}

\subsection{Held-out rewrite generalization}

To check whether the encoder is learning rewrite-rule signatures rather than circuit semantics, we retrain with only 4 of 5 rewrites and evaluate on positives generated by the held-out rewrite alone. We hold out \texttt{insert\_id} (HH/XX/CNOT-CNOT identity insertion), since this rewrite changes both gate counts and sorted gate lists and is therefore the most discriminating against baseline failure.

\begin{table}[h]
\centering
\begin{tabular}{lccc}
\toprule
method & R@1 & R@5 & R@10 \\
\midrule
hash & 0.000 & 0.000 & 0.000 \\
sorted hash & 0.000 & 0.000 & 0.000 \\
count vec & 0.045 & 0.160 & 0.290 \\
\textbf{quark (ours)} & \textbf{0.425} & \textbf{0.845} & \textbf{0.935} \\
\bottomrule
\end{tabular}
\caption{Held-out-rewrite retrieval. Model trained without \texttt{insert\_id}; eval pairs use \texttt{insert\_id} only.}
\end{table}

The 64.5-point R@10 advantage over count-vector at this distribution shift is direct evidence that the encoder learns transferable structural invariances --- in particular, it learned to look past redundant gate insertions even though it never saw \texttt{insert\_id} during training. We hypothesize that exposure to the inverse rewrite (\texttt{cancel\_consecutive}) during training induces this generalization.

\subsection{QASMBench: out-of-distribution retrieval}

We evaluate on the QASMBench small benchmark suite \cite{qasmbench} restricted to circuits with $\leq 8$ qubits and $\leq 128$ gates after parsing, yielding 29 algorithms each present in original and Qiskit-transpiled form. The retrieval task: given the original-form circuit, retrieve its transpiled lowering from a library of all 58 circuits.

\begin{table}[h]
\centering
\begin{tabular}{lccc}
\toprule
method & R@1 & R@5 & R@10 \\
\midrule
hash & 0.000 & 0.000 & 0.000 \\
sorted hash & 0.000 & 0.000 & 0.000 \\
count vec & 0.000 & 0.069 & 0.172 \\
\textbf{quark (ours)} & \textbf{0.000} & \textbf{0.241} & \textbf{0.345} \\
\bottomrule
\end{tabular}
\caption{QASMBench small: retrieve the transpiled form of an algorithm given its original-form circuit. 29 query/library pairs.}
\end{table}

This is a strongly out-of-distribution test: transpilation is not one of the rewrites used in training; transpiled circuits have substantially different gate counts and vocabularies (e.g., \texttt{adder\_n10}: 19 gates $\rightarrow$ 171 gates). Quark's 17.3-point R@10 lead over count-vector is real but absolute performance is weak. We do not claim Quark is yet a tool for cross-compilation circuit retrieval; the result is reported because hiding it would mislead.

\section{Limitations}

\begin{itemize}[topsep=2pt,itemsep=2pt]
    \item \textbf{Training distribution.} Random gate sequences with limited rewrite diversity. Real-world circuits (QFT variants, VQE ansatzes, error correction) likely require domain-specific finetuning.
    \item \textbf{Verifier.} \texttt{quark.verify.equiv} computes the full unitary; usable for ground truth on $\leq 10$ qubits, useless above.
    \item \textbf{No measurements.} The encoder vocabulary covers unitary gates only.
    \item \textbf{Equivalence as defined here is unitary equivalence up to global phase.} Equivalences arising from measurement-conditioned operations or relative-phase identities are out of scope.
    \item \textbf{QASMBench performance.} The +17 R@10 advantage suggests there is signal but the model is not solving the cross-compilation retrieval task in any practical sense.
\end{itemize}

\section{Future work}

Hard negative mining; finetuning on real-corpus distributions; expanded rewrite vocabulary including ZX-calculus identities and known compiler optimizations; measurement and classical-control support.

\section{Conclusion}

Quark is a small library and a small model. The contribution is structural: a clean baseline for circuit embedding, a generalization test that distinguishes structure-learning from rewrite-memorization, and an honest report of how the model performs on real-corpus retrieval. Code, weights, and benchmarks are reproducible from the public repository.

\bibliographystyle{plainnat}
\begin{thebibliography}{9}
\bibitem{qiskit} Qiskit Development Team. Qiskit: An Open-Source Framework for Quantum Computing. \url{https://www.ibm.com/quantum/qiskit}.
\bibitem{pennylane} V. Bergholm et al. PennyLane: Automatic differentiation of hybrid quantum-classical computations. arXiv:1811.04968, 2018.
\bibitem{infonce} A. v.d. Oord, Y. Li, O. Vinyals. Representation Learning with Contrastive Predictive Coding. arXiv:1807.03748, 2018.
\bibitem{clip} A. Radford et al. Learning Transferable Visual Models From Natural Language Supervision. ICML 2021.
\bibitem{qasmbench} A. Li, S. Stein, S. Krishnamoorthy, J. Ang. QASMBench: A Low-level QASM Benchmark Suite for NISQ Evaluation and Simulation. ACM TQC 2022.
\end{thebibliography}

\end{document}
